\section[Réévaluation]{Réévaluation}
\label{sec:etat-art-reevaluation}

Paragraphe : (Rapidement sur) les systèmes de modélisation géométriques pour la
CAO. (Principalement) les systèmes de modélisation paramétriques, (Concentrer
sur) la famille fonctionnelle, c'est-à--dire, avec spécification paramétrique.
(Introduire) les besoins de générations de modèles et (motiver rapidement avec) les coûts de changements dans les designs.

Dans un premier temps~: concepts, définitions et propriétés
Puis~: modélisation paramétrique (famille fonctionnelle)

\subsection[Concepts, définitions et propriétés]{Concepts, définitions et propriétés}
\label{sec:etat-art-reevaluation-defs-props}

\begin{definition}[entités invariantes]

\end{definition}

\begin{definition}[entités contingente]

\end{definition}

\begin{definition}[entités avec sources]

\end{definition}

\begin{definition}[entités sans sources]

\end{definition}

\begin{definition}[specification paramétrique]

\end{definition}

\begin{definition}[historique]

\end{definition}


\subsubsection[Modélisation paramétrique]{Modélisation paramétrique}
\label{sec:etat-art-reevaluation-model-param}

% La modélisation paramétrique une méthode de modélisation incontournable en
% CAO. %
% Dans ces systèmes, un utilisateur crée un objet 3D uniquement à partir
% d'opérations de modélisation à paramétrer. %
% Ces paramètres permettent de restituer au mieux les intentions de modélisation
% de l'utilisateur. %

\paragraph*{Modélisation paramétrique}

\paragraph{Spécification paramétrique}

\paragraph{Réévaluation}

\subsubsection[Nomination persistante : définitions et propriétés]{Nomination persistante : définitions et propriétés}
\label{sec:etat-art-reevaluation-model-param-nom-persistant-def-props}

\begin{definition}[nommage persistant]

\end{definition}

\subsubsection[Critères de classification des systèmes de modélisation]{Critères de classification des systèmes de modélisation}
\label{sec:etat-art-reevaluation-model-param-criteres}

\subsection[Approches existantes en CAO]{Approches existantes en CAO}
\label{sec:etat-art-reevaluation-approches}

\begin{itemize}
	\item Entités invariantes et contingentes
	\item Entités avec et sans source
	\item Nomination des entités issues d'extrusion ou de révolution
	\item Nomination des entités basée sur les entités invariantes
\end{itemize}

% \subsubsection[Distinction des entités invariantes et contingentes]{Distinction des entités invariantes et contingentes}
% \label{sec:etat-art-reevaluation-approches-dis-entites}

% \subsubsection[Distinction des entités avec et sans sources]{Distinctions des entités avec et sans sources}
% \label{sec:etat-art-reevaluation-approches-dis-sources}

% \subsubsection[Nomination des entités issues d'extrusion ou de révolution]{Nominations des entités issues d'extrusion ou de révolution}
% \label{sec:etat-art-reevaluation-approches-nomination-ext-rev}

% \subsubsection[Nomination basée sur les entités invariantes]{Nomination basée sur les entités invariantes}
% \label{sec:etat-art-reevaluation-approches-nomination-inv}

% \subsubsection[Utilisation d'un historique de noms]{Utilisation d'un historique de noms}
% \label{sec:etat-art-reevaluation-approche-historique-noms}

% \subsubsection[Nomination hétérogène des entités]{Nomination hétérogène des entités}
% \label{sec:etat-art-reevaluation-approche-nomination-heterogene}



\subsection[Réévaluation à base de règles]{Réévaluation à base de règles}
\label{sec:etat-art-reevaluation-reevaluation-regles}


%%% Local Variables:
%%% mode: latex
%%% TeX-master: "../../../main"
%%% End:
